\chapter[The Black Death and the Old World]{How the Black Death Changed the Old World Together}
\section{A Gold Coin from the Black Sea}
Pick up a fourteenth-century silver coin.

Small and thinner than a one-yuan coin, its irregular edges reflect hand-hammering: no two were identical. It came from Genoa's Black Sea trading station at Caffa in today's Crimea. Counted in hands, tucked in leather, carried deep in holds, it might have paid for grain, silk, or spices.

Imagine its route. The early fourteenth century offered exceptionally connected commerce. Mongol khanates linked Yuan Dadu to the Black Sea through relay stations. Historians call this the Pax Mongolica, referring not to Mongol kindness but unexpectedly secure trade. Chinese ceramics, Indian pepper, Central Asian jade, and Venetian glass moved with unprecedented ease through one network.

Goods were not its only passengers.

In 1346, Golden Horde troops besieging Caffa began dying horribly: fever, egg-sized swellings in armpits and groins, bleeding beneath skin, dark patches, and death within days. Genoese behind walls initially felt safe. Chronicler Gabriele de' Mussi reported corpses catapulted into the city. Widely repeated, the story is now generally considered unreliable; walls could not stop the real carriers anyway.

Black rats and their fleas hid in cargo. Fleas feeding on infected rats moved after their hosts died to another warm body, rat or human. In autumn 1347, Genoese ships escaping Caffa reached Messina in Sicily with many aboard sick or dead, bringing rats and fleas ashore.

Within days, Messina's people began dying.

Coin, siege, and ships connect our chapter: \textbf{an invisible pathogen followed trade and war to achieve what no emperor or army had done, bringing successive collapse from northern China to Icelandic fishing villages within little more than a decade}. Contemporaries called it the Great Pestilence or Great Mortality; we call it the Black Death. Medically, plague is caused by Yersinia pestis, named after French physician Alexandre Yersin, who isolated it during Hong Kong's 1894 epidemic.

Hold the silver coin and come along.

\section{A Florentine Morning in 1348}
Spring 1348: Florence, among Europe's richest cities, approaches a hundred thousand people. Wool and banking sustain it. English fleeces become cloth sold across half Europe; even France's king uses Florentine bankers' bills.

Ordinary mornings begin with bells, combers' chatter, and wool-cart wheels on paving. Now bells ring too often for the dead, until municipal authorities limit them because sound itself spreads panic.

Street smells mingle fragrance and stench. Burning herbs and spices answer medicine's miasma theory, disease carried by corrupted air. People wear scent balls and lift them to their noses. Unburied bodies supply the stench. Gravediggers are insufficient, cemeteries full, and trenches outside town stack corpses beneath thin soil layers.

People unsettle you more. Mothers avoid children; husbands fear wives' rooms. Remember: hearts did not suddenly turn wicked. Yesterday's dinner companion has swellings today. Nobody knows transmission's mechanism; everyone knows it \textbf{happens}. Fear rewrites customs: no embracing friends, vinegar bowls at shop entrances, coins accepted only after soaking.

The city struggles too, banning rubbish dumping, ordering dead people's clothing burned, questioning outsiders. Crude measures broadly approach hygiene and isolation through experience without microscopes.

Thirty-five-year-old Giovanni Boccaccio, from a wealthy family, witnesses it, including his stepmother's death. Years later, he opens the Decameron with ten young people fleeing to a country villa and telling a hundred stories, one each per day. We will soon read its famous introduction.

Similar mornings reach Paris, London, Cairo, Damascus, and Constantinople. A widely accepted estimate places Europe's 1347--1353 losses at roughly a third to a half: not one unfortunate village, an entire continent.

\section[A Flea Needs the World's Help]{A Flea Cannot Overturn the World Unless the World Supplies a Handle}
State the conflict before answers.

This appears biological: bacteria, fleas, rats, death. Stop there and plague resembles an earthquake, unpredictable disaster followed by survival, rebuilding, and history continuing.

But consider the subsequent decades:

English serfs refuse unpaid labour and demand rising wages. Parliamentary wage controls fail. Over thirty years later, tens of thousands revolt and nearly enter London to capture the king.

France experiences the Jacquerie peasant revolt in 1358.

Florentine wool combers, among Europe's earliest industrial workers, briefly seize political power in 1378.

Thousands in Germany and northern France march whipping themselves to cleanse humanity's sins; others carry torches into Jewish quarters.

Priests die in great numbers while administering last rites; survivors' preaching loses credibility. Over a century later, when Martin Luther posts his theses, many hearts already carry cracks opened by plague.

Yet equally devastated Cairo and Damascus experience no comparable transformation. Mamluk rule continues, peasants worsen, and eastern European landlords tighten control over subsequent centuries.

The real question appears:

\textbf{How can unintentional bacteria and fleas rewrite wages, law, faith, and power?}

If fewer people naturally change everything, why loosen serfdom in England but tighten it in Poland and Russia? If fear produces temporary chaos, why do wage and status changes last centuries? If divine punishment explains it, as millions believed, why do explanations themselves change afterwards?

The epidemic resembles an experiment introducing death into dozens of societies with different outcomes. How does disaster become history?

Before proceeding, wager intuitively: did mass death automatically transform Europe, or did survivors' responses do so?

\section{Survivors Calculate Differently}
An apparently peaceful surviving village contains competing calculations. Hear several voices.

\textbf{The surviving worker.} In spring 1349, a Kent serf notices his lord becoming not kinder but more \textbf{desperate}. A third of neighbours are dead; crops cannot await burial. Once he begged for land; now the lord begs for hands. Labour has value. Nearby manors secretly raise wages to recruit; why work free? Elite chronicles complain of arrogant servants refusing low pay and demanding double. Angry accounts from opponents provide strong evidence of improved bargaining power.

\textbf{The lord and king.} Strengthen their case, difficult though it is given the harm it later caused.

To a lord, manor life forms an inherited contract: peasants supply labour, lords protection and land, with customary rather than negotiated prices. Labour costs soar while contractual income stays fixed, threatening bankruptcy. Labour service is status, not a resignable job; fleeing for higher pay means broken obligation and social collapse. Kings calculate war costs too, amid the Hundred Years War. The 1351 Statute of Labourers represents an entire old order defending itself against price signals, not simple wickedness.

Only after hearing this can we see the problem: an agreement formed amid abundant labour is treated as eternal justice, and law commands markets to pretend plague never happened. We will see the result.

\textbf{The church and its cracks.} Priests anointing the dying face extreme danger. Heavy clerical mortality forces hurried ordination, sometimes after minimal training. Believers ask why masses fail and why priests die faster if their explanations are right.

Flagellants emerge through these cracks. In 1348--1349, barefoot groups cross Germany and Flanders in thirty-three-and-a-half-day cycles, matching Jesus' earthly years, whipping backs with metal-studded straps before weeping recruits. Before calling them foolish, consider their available explanation: invisible bacteria, one in three dead, sin angering God, and inability to do nothing. Volunteering to bear humanity's punishment has coherent, even tragic, logic. Understanding origins does not endorse direction. The movement later attacks clergy and Jews and is banned by the pope in 1349.

\textbf{The accused.} Now the chapter's heaviest passage.

Fear prefers one explanation: \textbf{someone poisoned us}. Rumours accused Jews of poisoning wells, with an unfalsifiable answer to every objection. Jewish deaths meant guilty concealment; deaths where no Jews lived meant contaminated water arriving elsewhere. In Strasbourg in February 1349, Jews were driven into their cemetery's execution ground and roughly two thousand killed, many having traded peacefully with Christian neighbours the previous day. Similar massacres spread across Mainz, Cologne, and dozens of cities. Tortured confessions in Savoy and elsewhere travelled by horse as evidence for further arrests.

Trace the chain: rumour, torture, false testimony, silencing victims. No stage requires real evidence.

Contemporary voices opposed it. Clement VI issued protective bulls in 1348, pointing out that plague killed Jews and struck places without them, refuting minority poisoning. Reason could not outrun torch-bearing crowds. Jewish communities nonetheless remained active: survivors recorded deaths in distant correspondence, buried relatives, and later rebuilt synagogues in looted streets. Disaster accounts must retain persecuted people's agency, not reduce them to scenic victims. They stitched life together again.

\textbf{Across the Mediterranean.} Cairo's losses rivalled Europe's. Seventeen-year-old Tunisian Ibn Khaldun lost both parents, later becoming a great historian. Instead of flagellation or poison accusations, he pursued study and eventually wrote the Muqaddimah about civilisations' rise and fall. Granadan physician Ibn al-Khatib defended real contagion despite religious objections, arguing that exposed people became ill while the unexposed stayed well. Humanity used divine, human, empirical, and political explanations, tested across the following century.

\section[Thinking Bridge: Four Disaster Layers]{Thinking Bridge: Divide an Epidemic into Four Layers}
Organise the material with a portable tool: \textbf{analyse any major disaster through biology, institutions, information, and morality}.

Disaster never operates on only one level. Separating them identifies distinct patterns and responsibilities; mixing them leaves only terror.

\textbf{Biology: what kills, and how does it move?}

Yersinia pestis travels mainly through rat fleas; crowded populations may experience pneumonic plague through respiratory droplets. Trade and grain ships carry it at ship rather than horse speed, striking ports first. Testable evidence now points towards the late-1330s outbreak near Kyrgyzstan's Lake Issyk-Kul, with inscriptions and ancient pathogen DNA matching place and time. Biology contains no villain: bacteria reproduce, not commit evil.

\textbf{Institutions: how do existing structures amplify or buffer shock?}

Who sets wages, controls land, commands law? English peasants could bargain among manors, converting scarcity into higher wages; eastern lords could prevent departure, converting it into stronger restraints. \textbf{Same shock, different institutions, different outcomes.} Amplification and buffering are human-made structures, not natural disaster.

\textbf{Information: what understanding guides action?}

What beliefs and evidence existed, and why did rumour outrun truth? Miasma, divine punishment, markets, and pulpits supplied scent balls, repentance, and suspicion at wells. Without microscopes, people could not know fleas' role; informational lag was not stupidity. But failures imposed costs, connecting to the next layer.

\textbf{Morality: whom did fear blame, and who paid?}

Why should Strasbourg's Jews pay with life for a disease killing Jews too? Why send priests into deadly rooms? Why criminalise serfs' bargaining? Each injustice involved concrete choices, not inevitability. Calling plague a shared human tragedy safely erases distinctions between misfortune and persecution. Record them separately.

The full picture appears: \textbf{biological migration, institutional stress test, explanatory failure across Europe, and a moral collapse into scapegoating}. Different rules at each level explain apparently contradictory results: the same illness meets different institutions and explanations.

Use this beyond medieval history, for epidemics, accidents, or viral news: what happened, which structures amplified it, what informed responses, who bore costs? Many quarrels merely shout between levels.

\section[Two Nearby Documents]{Reading Evidence: Two Documents Written Close Together}
Theory cannot replace surviving paper. Compare a writer's eyes and a lawmaker's seal.

\textbf{First: Boccaccio's Decameron introduction, around 1353.}

He records social breakdown with chilling composure: brothers abandon brothers, uncles nephews, wives husbands, even parents children. Another unforgettable detail describes a dead sick pig's rags---actually rags belonging to a dead person---tossed by two pigs, which then circle and collapse dead. His point is contagion so potent that touching the dead person's cloth kills.

Do not reduce this to declining morals. First, it is a \textbf{record of failed medical understanding}, not simply judgement. People saw affection transmit death without knowledge separating contact and danger. Love had not vanished; fear quarantined it. Second, the ensuing hundred stories freely mix bawdiness, laughter, and attacks on clergy and moralists. The book's structure reveals survivors' weakening reverence for old authority.

\textbf{Second: England's Statute of Labourers, enacted by Parliament in 1351.}

Its preamble openly resents workers who, after widespread mortality, exploit employers' need and their own scarcity by refusing old wages and demanding excessive pay, otherwise remaining idle. It fixes wages at pre-plague levels, punishes breaches with imprisonment, and penalises employers paying above the limits too.

Written only a few years apart, the documents seem planetary opposites: Boccaccio addresses morality and torn relationships; legislation tries institutionally to force the changed world back into shape. Repeated reinforcement and penalties reveal the law's failure: wages rise and serfs flee. In 1381, tens of thousands enter London demanding abolition of serfdom. A preaching rhyme attributed to John Ball asks:

\begin{quote}
When Adam dug and Eve wove, who was the nobleman?
\end{quote}
Without economics terminology, this dismantles fixed social places. If everyone began as equal workers, who decided when and why one was born lord and another serf? Plague did not invent the question; it gave questioners confidence.

Together these documents form our structure: biological disaster becomes both moral and institutional crisis within years.

\section{Fewer People, So Everything Changed?}
Separate events, contemporary understanding, and evaluation before comparing causes. Europe lost a third to a half, wages later rose, and serfdom loosened; people invoked God, air, and poison. Now ask \textbf{why} change followed.

\textbf{Explanation one: demographic determination---scarce labour changes everything.}

The chain is clear: mass death $\rightarrow$ fewer workers $\rightarrow$ scarcity raises prices $\rightarrow$ labour dearer, land cheaper $\rightarrow$ competing lords loosen control $\rightarrow$ serfdom dissolves, wages and ordinary lives improve. Some estimates roughly double English agricultural real wages over the following century; by the fifteenth, western serfdom was nominal and London's poor could afford meat and beer once marking status. Economics supplies powerful gears. Its limit comes next.

\textbf{Explanation two: institutions determine direction.}

This is our crucial \textbf{counterexample} to automatic improvement through population loss.

East of the Elbe, Poland, Prussia, and Russia also experienced mortality and scarcity, yet from the fifteenth century lords prohibited departure and increased labour obligations, producing second serfdom, harsher than the first. In Egypt, Mamluk military landholders raised survivors' taxes rather than wages, demanding payment even for the dead. Recovery slowed, depression persisted, and no workers' spring arrived.

One disease and scarcity yielded western release, eastern restriction, Egyptian extraction. \textbf{Population pushes; prior power and institutions determine where society falls.} Western peasants had mobility and legal room, lords competed, and kings welcomed weakened landlords. Eastern landlords monopolised politics amid fewer urban escape routes. This asks where social valves lie rather than predicting one demographic outcome. Yet institutional emphasis can understate mortality's own force: scarcity genuinely drove restrictive legislation.

\textbf{Explanation three: failed beliefs and psychological change.}

The church's former explanatory monopoly cracked when prayer failed, priests died, and papal decrees could not stop torches. Space opened for physicians such as Ibn al-Khatib, writers such as Boccaccio, and preachers such as Ball. Some trace Reformation and Scientific Revolution roots here. This addresses changing minds beyond wages, but nearly 170 years separate 1348 and Luther's 1517, with disputed intermediate links. Treat it as loosened soil, not a straight causal arrow.

\textbf{The first explains motive force, the second direction, the third echoes}. Different ranges complement rather than exclude. Keep the counterexample: anyone claiming one factor necessarily produces one outcome must explain eastern Europe's opposite path.

\section[Further Challenge: Germs in History]{Further Challenge: Invite Germs Back into History}
Why did disasters on this scale scarcely feature in histories before the fourteenth century? Our theory addresses that broader question.

Traditional history favours rulers, wars, treaties, and ideas: human intention. Here the greatest actor had none. William McNeill's 1976 Plagues and Peoples controversially restored pathogens to history's main line as variables in civilisation's rise and fall, helping establish historical epidemiology.

Two concepts are especially useful.

First, \textbf{microparasitism and macroparasitism}. Parasites live off others. Bacteria and viruses extract from humans through mortality; armies, tax collectors, and lords through taxes and labour. The harsh analogy asks whether either exhausts its host. Their combination invites disaster: 1340s Europe faced war, taxes, and recent famine before plague. Microparasites undermined macroparasites by collapsing populations, tax bases, and old channels of extraction.

Second, \textbf{virgin-soil epidemics}: pathogens entering unexposed populations without immunity or experience can produce catastrophic mortality. The Black Death was such an epidemic for the fourteenth-century Old World. McNeill maps Eurasian peoples' long coexistence with local pathogens, immunity purchased through generations of death rather than innate strength. Mongol-era connections meant \textbf{goods networks also connected pathogen reservoirs}. Messina was not first: sixth-century Justinianic plague followed grain routes; Justinian recovered, the empire did not. Nor last: European ships brought smallpox to the Americas, with devastating losses among unprepared Indigenous peoples. Totals remain disputed, but an epidemic-led collapse is widely accepted; Aztec and Inca courage could not compensate for immune vulnerability.

This enlarges civilisational encounters from guns and ships to immunological histories, changing how we ask why the Old World conquered the New rather than vice versa.

Yet its limits match our four layers. McNeill explains pathogen arrival and vulnerability, \textbf{not Strasbourg's fires}. Bacteria ordered no torches; pathogen reservoirs contain no account of scapegoats. Biology must hand over to institutions, information, and morality. Knowing a theory's shadow matters as much as knowing its light.

\section[Back to Today: The Pandemic Years]{Back to Today: The Three Years We Just Experienced}
Between writer and reader stands shared experience of the COVID-19 pandemic beginning in 2020. Apply our four layers.

\textbf{Biology: very different numbers.} The Black Death killed a third to a half of Europe; WHO estimated about fifteen million pandemic-associated excess deaths in 2020--2021, enormous but roughly two per thousand of an eight-billion world. The difference reflects four centuries of biomedical accumulation: identifying pathogens, testing, vaccines, and rescuing severe patients with ventilators and antibiotics. Establish this before comparison.

\textbf{Institutions: familiar patterns.} Quarantine's name comes from this legacy. In 1377 Ragusa, today's Dubrovnik, required incoming people and ships to wait thirty days outside port; later Venice and others used forty, quaranta in Italian. Six centuries later came lockdowns, temporary hospitals, health codes, online classes, and disrupted supply chains. Different systems combined restrictions differently, with different costs and benefits, as labour scarcity once divided European serfdom's fortunes. \textbf{Disasters do not create institutional differences; they apply high voltage to existing ones.}

\textbf{Information: almost a high-definition replay of 1348.} Then poisoned-well rumours; recently claims of engineered viruses, vaccine chips, or universal-cure oral remedies. Papal decrees lost to torches; public-health rebuttals lose to family-chat forwards. \textbf{Fear needs one story; evidence needs verification, so fear runs faster}. Progress exists too: Ibn al-Khatib risked religious accusation for contagion, whereas public discussion now permits immediate peer challenge.

\textbf{Morality: closest to us.} Did hearing of a returning traveller first prompt resentment rather than concern for observation? Did you see barred delivery workers or avoided recovered patients? Reducing someone to a walking virus resembles the mechanism behind Strasbourg's gaze, despite vastly different scale. Fear seeks visible targets among marginal, foreign, less-heard people. Our advantage is not nobler nature but more monitoring, faster correction, and laws constraining exclusion. Neither despair nor complacency follows: \textbf{scapegoating survives in a cage whose key each generation holds}.

Wages echo too. Post-pandemic labour shortages, pay increases, and Great Resignation debates revived worker leverage. Is this the logic preceding the 1351 statute? Scholars debate; direct equivalence is premature. The question shows the Black Death's continued historical presence in changed clothing.

\section[Your Question: What Comes First?]{Your Question: What Should We Save First in the Next Epidemic?}
Return to the silver coin.

Its world knew neither bacteria nor flea control; rumours outran ships and fear burned neighbourhoods. Institutions turned one disease towards loosened or tightened chains. Today, microscopes, vaccines, and internet access permit something survivors could scarcely imagine: discussing priorities \textbf{in advance}.

Take the wheel. A highly contagious new disease appears, mortality unknown, information confused. Several needs become urgent but resources can prioritise only one or two:

Protect \textbf{truth}, publishing all data including frightening and unverified bad news, or \textbf{order}, delaying and tiering release to steady people despite lost opportunities for self-protection? Protect the \textbf{economy}, keeping workplaces open despite illness, or the \textbf{most vulnerable}, imposing shutdowns to protect elders, patients, and those too poor to stop working despite unemployment and hunger? Protect \textbf{freedom} through voluntary isolation, or the \textbf{collective} through compulsory quarantine like Ragusa's forty days, risking unnecessary confinement rather than missed cases?

Rank priorities and justify them through consequences and acceptable or unacceptable costs, not mere preference.

Fourteenth-century survivors had no opportunity to discuss beforehand; they answered weeping beside bodies. You have that opportunity. When the day comes, you need not begin learning fear from nothing.
